\documentclass[a4paper,12pt]{article}
\usepackage[utf8]{inputenc}
\usepackage[english,russian]{babel}
\usepackage[T2A]{fontenc}
\usepackage[
  a4paper, mag=1000, includefoot,
  left=1.1cm, right=1.1cm, top=1.2cm, bottom=1.2cm, headsep=0.8cm, footskip=0.8cm
]{geometry}

\usepackage{amsmath}
\usepackage{amssymb}
\usepackage{times}
\usepackage{mathptmx}
\usepackage{listings}
\usepackage{xcolor}

\definecolor{codegray}{rgb}{0.5,0.5,0.5}
\definecolor{backcolour}{rgb}{0.95,0.95,0.92}

\IfFileExists{pscyr.sty}
{
  \usepackage{pscyr}
  \def\rmdefault{ftm}
  \def\sfdefault{ftx}
  \def\ttdefault{fer}
  \DeclareMathAlphabet{\mathbf}{OT1}{ftm}{bx}{it} % bx/it or bx/m
}

\mathsurround=0.1em
\clubpenalty=1000%
\widowpenalty=1000%
\brokenpenalty=2000%
\frenchspacing%
\tolerance=2500%
\hbadness=1500%
\vbadness=1500%
\doublehyphendemerits=50000%
\finalhyphendemerits=25000%
\adjdemerits=50000%


\begin{document}
\renewcommand{\theequation}{1.\arabic{equation}}
\author{Сергей Криворученко}
\title{Метод Холецкого решения линейной системы с симметричной матрицей}
\date{\today}
\maketitle

\begin{center}
{\bfseries \Large\underline{Постановка задачи}}
\end{center}

Дана симметричная матрица $A$ и вектор $b$ ищем решение системы $Ax = b$

$$
A=
  \begin{pmatrix}
    a_{11}& a_{12} &\ldots & a_{1n} \\
    a_{12}& a_{22} &\ldots & a_{2n} \\
    \vdots& \vdots &\ddots & \vdots\\
    a_{1n}& a_{2n} &\ldots & a_{nn}
  \end{pmatrix}
\qquad
b=
  \begin{pmatrix}
    b_{1} \\
    b_{2} \\
    \vdots \\
    b_{n}
  \end{pmatrix}
$$
Представим матрицу в виде:
$$
\begin{gathered}
A=R^T D R, \\
\text{где } R=
\begin{pmatrix}
    R_{11}^{m \times m} & R_{12}^{m \times m} & \cdots & R_{1k}^{m \times m} & R_{1k+1}^{m \times l} \\
    0 & R_{22}^{m \times m} & \cdots & R_{2k}^{m \times m} & R_{2k+1}^{m \times l} \\
    \vdots & \vdots & \ddots & \vdots & \vdots \\
    0 & 0 & \cdots & R_{kk}^{m \times m} & R_{kk+1}^{m \times l} \\
    0 & 0 & \cdots & 0 & R_{k+1k+1}^{l \times l}
\end{pmatrix},
\qquad
D=
\begin{pmatrix}
  D_{11}^{m \times m} & 0 & \cdots & 0 & 0 \\
  0 & D_{22}^{m \times m} & \cdots & 0 & 0 \\
  \vdots & \vdots & \ddots & \vdots & \vdots \\
  0 & 0 & \cdots & D_{kk}^{m \times m} & 0 \\
  0 & 0 & \cdots & 0 & D_{k+1k+1}^{l \times l}
\end{pmatrix} \\
\text{где $R_{ii}$ - верхнетреугольная матрица, а $D_{ii}$ - диагональная матрица с $\pm1$ на диагонали}
\end{gathered}
$$

$m$ - заданный размер блока матрицы, а $l$ - остаток от деления $n$ на $m$

\section{Хранение}

\subsection{Способы хранения}

Хранить матрицы будем одномерными массивами длины
$k \times l$, $\frac{n(n+1)}{2}$ и $n$

\medskip

Матрица $A$ хранится по строкам. В силу симметрии сохраняется только
верхний треугольник, поэтому размер массива равен
$$
\frac{n(n+1)}{2}
$$

Блок $A_{ij}$ также хранится по строкам, но целиком, в массиве размера
$o \times p$, где
$$
o =
\begin{cases}
  l, & i = k+1, \\
  m, & i \leq k
\end{cases}
\qquad
p =
\begin{cases}
  l, & j = k+1, \\
  m, & j \leq k
\end{cases}
$$

\medskip

Матрица $R$ хранится по столбцам. Это можно рассматривать как
хранение $R^T$ по строкам. Сохраняется только верхний треугольник
в массиве размера
$$
\frac{n(n+1)}{2}
$$

Блок $R_{ij}$, аналогично $A_{ij}$, хранится целиком даже при
$i=j$, поскольку экономия памяти в данном случае не оправдывает
дополнительные вычисления. Элементы блока записываются по столбцам

\medskip

Диагональная матрица $D$ хранится в массиве типа \texttt{int}
размера $n$. В массив записываются только диагональные элементы

\medskip

Доступ к элементам зависит от порядка хранения и размера массива

\medskip
\noindent
\textbf{Массив размера $o \times p$}

Для $1 \leq i \leq o$, $1 \leq j \leq p$:
$$
\begin{aligned}
  \text{по строкам:}\qquad
  a[(i-1)p+j-1] &= a_{ij}, \\
  \text{по столбцам:}\qquad
  a[(j-1)o+i-1] &= a_{ij}
\end{aligned}
$$

\medskip
\noindent
\textbf{Массив размера $\frac{n(n+1)}{2}$}

Для $1 \leq i \leq j \leq n$:
$$
\begin{aligned}
  \text{по строкам:}\qquad
  a\left[(i-1)n-\frac{i(i-1)}{2}+j-1\right]
  &= a_{ij}, \\
  \text{по столбцам:}\qquad
  a\left[\frac{j(j-1)}{2}+i-1\right]
  &= a_{ij}
\end{aligned}
$$

\medskip
\noindent
\textbf{Массив размера $n$}

Для $1 \leq i \leq n$:
$$
a[i-1] = a_{ii}
$$

\subsection{Чтение / Запись блоков}

Для операций с блоками будем вызывать функции

void get\_block\_row (const double* a, double *b, const int n, const int m, const int l, int si, int sj);

void set\_block\_row (double* a, const double *b, const int n, const int m, const int l, int si, int sj);

void get\_block\_column (const double* a, double *b, const int n, const int m, const int l, int si, int sj);

void set\_block\_column (double* a, const double *b, const int n, const int m, const int l, int si, int sj);

a - указатель на матрицу

b - указатель на массив для блока

si, sj - указатели на начало блока

n - размер матрицы

m - высота блока

l - ширина блока

\begin{lstlisting}[language=C++,
                   tabsize=4,
                   stepnumber=1,
                   breaklines=true,
                   numbers=left,
                   backgroundcolor=\color{backcolour},
                   basicstyle=\ttfamily\small,
                   columns=fullflexible,
                   keepspaces=true,
                   numbersep=8pt,
                   lineskip=1pt
                  ]
void get_block_row (const double *a, double *b, const int n, const int m, const int l, int si, int sj)
{
  if (si < sj)
    for (int i = 0; i < m; ++i)
      {
        int row = si + i;
        
        int pos = row * n - (row * (row + 1) >> 1) + sj;
        
        std::memcpy (b + i * l, a + pos, l * sizeof(double));
      }
  else if (si == sj)
    for (int i = 0; i < m; ++i)
      {
        int row = si + i;
        int pos = row * n - (row * (row + 1) >> 1) + row;
        int len = m - i;
        
        memcpy(b + i * m + i, a + pos, len * sizeof(double));
        
        for (int j = 0; j < i; ++j)
          b[i * m + j] = b[j * m + i];
      }
}

void set_block_row (double *a, const double *b, const int n, const int m, const int l, int si, int sj)
{
  if (si < sj)
    for (int i = 0; i < m; ++i)
      {
        int row = si + i;
        
        int pos = row * n - (row * (row + 1) >> 1) + sj;
        
        std::memcpy (a + pos, b + i * l, l * sizeof(double));
      }
  else if (si == sj)
    for (int i = 0; i < m; ++i)
      {
        int row = si + i;

        int pos = row * n - (row * (row + 1) >> 1) + row;

        int len = m - i;
        
        memcpy(a + pos, b + i * l + i, len * sizeof(double));
      }
}

void get_block_column (const double *a, double *b, const int n, const int m, const int l, int si, int sj)
{
  if (si < sj)
    for (int j = 0; j < l; ++j)
      {
        int column = sj + j;
        
        int pos = (column * (column + 1) >> 1) + si;
        
        std::memcpy (b + j * m, a + pos, m * sizeof(double));
      }
  else if (si == sj)
    for (int j = 0; j < l; ++j)
      {
        int column = sj + j;
        
        int pos = (column * (column + 1) >> 1) + si;
        
        memcpy(b + j * m, a + pos, (j + 1) * sizeof(double));
        memset(b + j * m + j + 1, 0, (l - j - 1) * sizeof(double));
      }
}

void set_block_column (double *a, const double *b, const int n, const int m, const int l, int si, int sj)
{
  if (si < sj)
    for (int j = 0; j < l; ++j)
      {
        int column = sj + j;
        
        int pos = (column * (column + 1) >> 1) + si;
        
        std::memcpy (a + pos, b + j * m, m * sizeof(double));
      }
  else if (si == sj)
    for (int j = 0; j < l; ++j)
      {
        int column = sj + j;
        
        int pos = (column * (column + 1) >> 1) + si;
        
        memcpy(a + pos, b + j * m, (j + 1) * sizeof(double));
      }
}
\end{lstlisting}

\section{Алгоритм}

\subsection{Преобразование матрицы}

Представим матрицу $A$ в виде

$$
A = R^T D R
$$

Поделим матрицу на $(k+1) \times (k+1)$ блоков

Рассмотрим первый диагональный блок:

$$
A_{11} = R_{11}^T D_{11} R_{11}
$$


\subsubsection{Обычный метод Холецкого}

Матрица $A_{11}^{m \times m}$ симметрична, значит её можно представить
в аналогичном виде

Для $1 \leq i \leq j \leq m$:

$$
\left[
\begin{aligned}
  d_{ii} &= \operatorname{sign}(a_{ii} - \sum_{k=1}^{i-1} r_{ki}^2 d_{kk}) \\
  r_{ii} &= \sqrt{\left| a_{ii} - \sum_{k=1}^{i-1} r_{ki}^2 d_{kk} \right|} \\
  r_{ij} &= \frac{a_{ij} - \sum_{k=1}^{i-1} r_{ki} r_{kj} d_{kk}}{r_{ii} d_{ii}}
\end{aligned}
\right.
$$

Таким образом найдём матрицы $R_{11}$ и $D_{11}$


\subsubsection{Поиск блоков}

Рассмотрим блоки первой блочной строки

Для $2 \leq j \leq k$:

$$
A_{1j}^{m \times m} = R_{11}^T D_{11} R_{1j}
$$

Для $1 \leq i, s \leq m$:

$$
(A_{1j})_{is}
=
\sum_{t=1}^{i-1}
(R_{11})_{ti} (R_{1j})_{ts} d_{tt}
+
(R_{11})_{ii} (R_{1j})_{is} d_{ii}
$$

Отсюда

$$
(R_{1j})_{is}
=
\frac{
  (A_{1j})_{is}
  -
  \sum_{t=1}^{i-1}
  (R_{11})_{ti} (R_{1j})_{ts} d_{tt}
}{
  (R_{11})_{ii} d_{ii}
}.
$$

Для $R_{1k+1}^{m \times l}$ вычисления аналогичны. Отличается только
размерность матрицы:

$$
A_{1k+1}^{m \times l}
=
(R_{11}^T)^{m \times m}
\cdot
D_{11}^{m \times m}
\cdot
R_{1k+1}^{m \times l}.
$$

Таким образом получим

$$
R_{11}^{m \times m},
\ldots,
R_{1k}^{m \times m}
\text{ и }
R_{1k+1}^{m \times l}.
$$


\medskip

\textbf{Общий случай}

Далее мы движемся построчно сверху-вниз, вычисляя каждую строку слева-направо шаг за шагом, аналогично предыдущему блоку

Рассмотрим $R_{ij}$, где $1 \leq i \leq j \leq k+1$

$$
A^{n \times n}
=
(R^T)^{n \times n}
\cdot
D^{n \times n}
\cdot
R^{n \times n}
$$

$$
\begin{aligned}
  \text{где } A &=
  \begin{pmatrix}
    A_{11}^{m \times m}& A_{12}^{m \times m} &\ldots & A_{1i}^{m \times m} &\ldots & \vline & A_{1j}^{m \times m} & \vline &\ldots & A_{1k+1}^{m \times l} \\
    (A_{12}^{m \times m})^T& A_{22}^{m \times m} &\ldots & A_{2i}^{m \times m} &\ldots & \vline & A_{2j}^{m \times m} & \vline &\ldots & A_{2k+1}^{m \times l} \\
    \vdots& \vdots &\ddots & \vdots &\ddots & \vline & \vdots & \vline &\ddots & \vdots \\
    \hline
    (A_{1i}^{m \times m})^T& (A_{2i}^{m \times m})^T &\ldots & A_{ii}^{m \times m} &\ldots & \vline & A_{ij}^{m \times m} & \vline &\ldots & A_{ik+1}^{m \times l} \\
    \hline
    \vdots& \vdots &\ddots & \vdots &\ddots & \vline & \vdots & \vline &\ddots & \vdots \\
    (A_{1j}^{m \times m})^T& (A_{2j}^{m \times m})^T &\ldots & (A_{ij}^{m \times m})^T &\ldots & \vline & A_{jj}^{m \times m} & \vline &\ldots & A_{jk+1}^{m \times l} \\
    \vdots& \vdots &\ddots & \vdots &\ddots & \vline & \vdots & \vline &\ddots & \vdots \\
    (A_{1k+1}^{m \times l})^T& (A_{2k+1}^{m \times l})^T &\ldots & (A_{ik+1}^{m \times l})^T &\ldots & \vline & (A_{jk+1}^{m \times l})^T & \vline &\ldots & A_{k+1k+1}^{l \times l}
  \end{pmatrix} \\
  R^T &=
  \begin{pmatrix}
    (R_{11}^{m \times m})^T & 0 & \cdots & 0 & \cdots & 0 \\
    (R_{12}^{m \times m})^T & (R_{22}^{m \times m})^T & \cdots & 0 & \cdots & 0 \\
    \vdots & \vdots & \ddots & \vdots & \ddots & \vdots \\
    \hline
    (R_{1i}^{m \times m})^T & (R_{2i}^{m \times m})^T & \cdots & (R_{ii}^{m \times m})^T & \cdots & 0 \\
    \hline
    \vdots & \vdots & \ddots & \vdots & \ddots & \vdots \\
    (R_{1k+1}^T)^{l \times m} & (R_{2k+1}^T)^{l \times m} & \cdots & (R_{ik+1}^T)^{l \times m} & \cdots & (R_{k+1k+1}^T)^{l \times l}
  \end{pmatrix} \\
  D \cdot R &=
  \begin{pmatrix}
    D_{11}^{m \times m} R_{11}^{m \times m} & D_{11}^{m \times m} R_{12}^{m \times m} & \cdots & \vline & D_{11}^{m \times m} R_{1j}^{m \times m} & \vline & \cdots & D_{11}^{m \times m} R_{1k+1}^{m \times l} \\
    0 & D_{22}^{m \times m} R_{22}^{m \times m} & \cdots & \vline & D_{22}^{m \times m} R_{2j}^{m \times m} & \vline & \cdots & D_{22}^{m \times m} R_{2k+1}^{m \times l} \\
    \vdots & \vdots & \ddots & \vline & \vdots & \vline & \ddots & \vdots \\
    0 & 0 & \cdots & \vline & D_{jj}^{m \times m} R_{jj}^{m \times m} & \vline & \cdots & D_{jj}^{m \times m} R_{jk+1}^{m \times l} \\
    \vdots & \vdots & \ddots & \vline & \vdots & \vline & \ddots & \vdots \\
    0 & 0 & \cdots & \vline & 0 & \vline & \cdots & D_{k+1k+1}^{l \times l} R_{k+1k+1}^{l \times l}
  \end{pmatrix}
\end{aligned}
$$

Из произведения блоков следует:

$$
A_{ij}^{o \times p}
=
\sum_{t=1}^{\operatorname{min}(i, j)}
(R_{ti}^T)^{o \times o_t}
\cdot
D_{tt}^{o_t \times o_t}
\cdot
R_{tj}^{o_t \times p}
$$

Поскольку $i \leq j$, выделим слагаемое при $t=i$:

$$
R_{ii}^T \cdot D_{ii} \cdot R_{ij}
=
A_{ij}
-
\sum_{t=1}^{i-1}
R_{ti}^T \cdot D_{tt} \cdot R_{tj}
=
Y
$$

Произведения

$$
R_{ti}^T \cdot D_{tt} \cdot R_{tj}
$$

к этому моменту уже вычислены на предыдущих шагах.

Поэтому ещё при получении $R_{tj}$ можно сразу вычитать произведения из соответствующих нижних блоков матрицы $A$. Чуть ниже мы рассмотрим это поподробнее


В результате для поиска $R_{ij}$ необходимо решить

$$
R_{ii}^T \cdot D_{ii} \cdot R_{ij} = Y
$$

Матрицу $R_{ij}$ можно найти по той же схеме, что и $R_{1j}$:

$$
(R_{ij})_{rs}
=
\frac{
  (Y_{ij})_{rs}
  -
  \sum_{t=1}^{r-1}
  (R_{ii})_{tr} (R_{ij})_{ts} d_{tt}
}{
  (R_{ii})_{rr} d_{rr}
}
$$

После получения блока $R_{ij}$, нужно вычесть его из нижних зависимых блоков, для этого последовательно рассматриваем все блоки
$A_{rj}$, где
$$
i < r \leq j
$$

Из каждого блока $A_{rj}$ вычитаем произведение
$$
R_{ir}^{T} \cdot D_{ii} \cdot R_{ij},
$$
то есть выполняем обновление
$$
A_{rj}
\leftarrow
A_{rj}
-
R_{ir}^{T} \cdot D_{ii} \cdot R_{ij}
$$

Произведение за счёт диагональности $D$ вычисляем таким образом:

$$
(R_{ir}^{T} \cdot D_{ii} \cdot R_{ij})_{pq} = \sum_{t = 1}^{m} (R_{ir})_{tp} (R_{ij})_{tq} (D_{ii})_{tt} 
$$

При $i=j$ матрицы $R_{ii}$ и $D_{ii}$ вычисляются методом Холецкого,
описанным выше

Вычитание в данном случае не требуется, так как ниже $R_{ii}$ нет ненулевых матриц

Таким образом последовательно вычисляются блоки матриц $R$ и $D$

После завершения алгоритма получаем искомое разложение
$$
A = R^{T} \cdot D \cdot R.
$$

\subsection{Поиск решения}

Имеем систему линейных уравнений

$$
(R^T \cdot D \cdot R)x = b
$$

Разобьём решение системы на три шага:

\begin{enumerate}
  \item $R^T z = b$;
  \item $D y = z$;
  \item $R x = y$
\end{enumerate}


\subsubsection{Шаг 1}

Решим систему

$$
R^T z = b.
$$

В блочном виде:

$$
\begin{pmatrix}
  (R_{11}^{m \times m})^T & 0 & \cdots & 0 & \cdots & 0 \\
  (R_{12}^{m \times m})^T & (R_{22}^{m \times m})^T & \cdots & 0 & \cdots & 0 \\
  \vdots & \vdots & \ddots & \vdots & \ddots & \vdots \\
  \hline
  (R_{1i}^{m \times m})^T & (R_{2i}^{m \times m})^T & \cdots & (R_{ii}^{m \times m})^T & \cdots & 0 \\
  \hline
  \vdots & \vdots & \ddots & \vdots & \ddots & \vdots \\
  (R_{1k+1}^T)^{l \times m} & (R_{2k+1}^T)^{l \times m} & \cdots & (R_{ik+1}^T)^{l \times m} & \cdots & (R_{k+1k+1}^T)^{l \times l}
\end{pmatrix}
\cdot
\begin{pmatrix}
  z_1 \\
  z_2 \\
  \vdots \\
  z_i \\
  \vdots \\
  z_{k+1}
\end{pmatrix}
=
\begin{pmatrix}
  b_1 \\
  b_2 \\
  \vdots \\
  b_i \\
  \vdots \\
  b_{k+1}
\end{pmatrix}
$$

Сначала рассматриваем первый блок:

$$
R_{11}^T z_1 = b_1,
$$

где $R_{11}^T$ является нижнетреугольной матрицей

Для первого элемента:

$$
(z_1)_1
=
\frac{(b_1)_1}{(R_{11}^T)_{11}}
$$

Для произвольного элемента:

$$
(b_1)_i
=
\sum_{t=1}^{i}
(R_{11}^T)_{it}(z_1)_t
$$

Последовательно вычитая уже известные слагаемые, получаем

$$
(z_1)_i
=
\frac{
  (b_1)_i
  -
  \sum_{t=1}^{i-1} (R_{11}^T)_{it} (z_1)_t
}{
  (R_{11}^T)_{ii}
}.
$$

В результате получим вектор $z_1^m$

\medskip

Теперь рассмотрим общий случай. При вычислении сверху вниз имеем

$$
\sum_{t=1}^{i} R_{ti}^T z_t = b_i
$$

Произведения $R_{ti}^T z_t$ для $1 \leq t \leq i-1$ уже известны,
поэтому остаётся решить

$$
R_{ii}^T z_i
=
b_i
-
\sum_{t=1}^{i-1} R_{ti}^T z_t.
$$

Это система с нижнетреугольной матрицей $R_{ii}^T$, поэтому она
решается аналогично случаю $R_{11}^T z_1=b_1$

Таким образом последовательно получим

$$
z_1^m,\ldots,z_k^m,z_{k+1}^l,
$$

а значит и весь вектор $z^n$


\subsubsection{Шаг 2}

Решаем

$$
Dy=z
$$

Так как $D$ диагональная матрица с элементами $\pm1$, вычисление
выполняется непосредственно:

$$
y_i = D_{ii} \cdot z_i
$$


\subsubsection{Шаг 3}

Решаем систему

$$
Rx=y
$$

В блочном виде:

$$
\begin{pmatrix}
  R_{11}^{m \times m} & R_{12}^{m \times m} & \cdots & R_{1i}^{m \times m} & \cdots & R_{1k+1}^{m \times l} \\
  0 & R_{22}^{m \times m} & \cdots & R_{2i}^{m \times m} & \cdots & R_{2k+1}^{m \times l} \\
  \vdots & \vdots & \ddots & \vdots & \ddots & \vdots \\
  \hline
  0 & 0 & \cdots & R_{ii}^{m \times m} & \cdots & R_{ik+1}^{m \times l} \\
  \hline
  \vdots & \vdots & \ddots & \vdots & \ddots & \vdots \\
  0 & 0 & \cdots & 0 & \cdots & R_{k+1k+1}^{l \times l}
\end{pmatrix}
\cdot
\begin{pmatrix}
  x_1 \\
  x_2 \\
  \vdots \\
  x_i \\
  \vdots \\
  x_{k+1}
\end{pmatrix}
=
\begin{pmatrix}
  y_1 \\
  y_2 \\
  \vdots \\
  y_i \\
  \vdots \\
  y_{k+1}
\end{pmatrix}
$$

Вычисления аналогичны шагу 1. Отличие состоит в направлении обхода:
поскольку $R$ верхнетреугольная, неизвестные вычисляются снизу вверх

В результате получим вектор $x^n$, являющийся решением исходной
системы

\section{Оценка сложности}

\begin{enumerate}
  \item Количество операций на нахождение матриц $R_{11}$ и $D_{11}$ по формулам, приведенным в первом шаге

  Для нахождения суммы вида $\sum_{s = 1}^{i - 1} a b$ нужно $(i - 1) + (i - 2)$ действий

  Таким образом для вычисления $r_{ij}$ и $r_{ii}$ требуется $2i - 2$ операций

  А всего нужно

  $$\sum_{i = 1}^{m} (m + 1 - i) (2i - 2) = 2\sum_{t = 0}^{m - 1} (mt + t^2) = m^2 (m - 1) + \frac{2m^3 - 3m^2 + m}{6} = \frac{4}{3} m^3 - \frac{3}{2} m^2 + \frac{1}{6} m$$
\end{enumerate}

Будем считать, что

$$
n=km,
$$

то есть матрица полностью разбивается на \(k\times k\) блоков размера \(m\times m\)

Все арифметические операции считаем равными по стоимости. Умножение на элементы матрицы \(D\), равные \(\pm1\), за отдельную операцию не считаем

\subsection{Нахождение матриц \(R\) и \(D\)}

\subsubsection{Диагональные блоки \(R_{ii}\) и \(D_{ii}\)}

Для диагонального элемента используем формулу

$$
r_{ii}
=
\sqrt{
  \left|
  a_{ii}
  -
  \sum_{t=1}^{i-1}r_{ti}^{2}d_{tt}
  \right|
}
$$

Для вычисления суммы требуется \(i-1\) умножений и \(i-1\) вычитаний. После этого выполняется одна операция извлечения квадратного корня

Следовательно, для вычисления \(r_{ii}\) требуется

$$
2(i-1)+1=2i-1
$$

операций

Для внедиагонального элемента

$$
r_{ij}
=
\frac{
  a_{ij}
  -
  \sum_{t=1}^{i-1}r_{ti}r_{tj}d_{tt}
}{
  r_{ii}d_{ii}
}
$$

аналогично требуется \(i-1\) умножений, \(i-1\) вычитаний и одно деление, то есть

$$
2i-1
$$

операций

В \(j\)-м столбце верхнетреугольной матрицы \(R_{ii}\) вычисляются элементы с номерами

$$
i=1,\ldots,j
$$

Поэтому количество операций для одного диагонального блока равно

$$
\sum_{j=1}^{m}\sum_{i=1}^{j}(2i-1)
$$

Так как

$$
\sum_{i=1}^{j}(2i-1)=j^2,
$$

получаем

$$
\sum_{j=1}^{m}j^2
=
\frac{m(m+1)(2m+1)}{6}
$$

Диагональных блоков всего \(k\), поэтому

$$
C_1
=
\frac{k\,m(m+1)(2m+1)}{6}
$$

\subsubsection{Внедиагональные блоки \(R_{ij}\)}

Для \(i<j\) блок \(R_{ij}\) находится из уравнения

$$
R_{ii}^{T}D_{ii}R_{ij}=A_{ij},
$$

где из \(A_{ij}\) уже были вычтены вклады предыдущих блочных строк

Элемент блока вычисляется по формуле

$$
(R_{ij})_{rs}
=
\frac{
  (A_{ij})_{rs}
  -
  \sum_{t=1}^{r-1}
  (R_{ii})_{tr}(R_{ij})_{ts}d_{tt}
}{
  (R_{ii})_{rr}d_{rr}
}
$$

Для одного элемента требуется \(r-1\) умножений, \(r-1\) вычитаний и одно деление:

$$
2(r-1)+1=2r-1
$$

В каждой строке находится \(m\) элементов. Следовательно, для одного блока \(R_{ij}^{m\times m}\)

$$
C_{R_{ij}}
=
m\sum_{r=1}^{m}(2r-1)
$$

Поскольку

$$
\sum_{r=1}^{m}(2r-1)=m^2,
$$

имеем

$$
C_{R_{ij}}=m^3
$$

Число внедиагональных блоков верхнетреугольной матрицы равно

$$
\frac{k(k-1)}{2}
$$

Следовательно,

$$
C_2
=
m^3\frac{k(k-1)}{2}
$$

\subsubsection{Вычитание уже вычисленных произведений}

После вычисления блока \(R_{ij}\) выполняется обновление блоков

$$
A_{rj}
\leftarrow
A_{rj}
-
R_{ir}^{T}D_{ii}R_{ij}
$$

Для одного элемента произведения имеем

$$
(R_{ir}^{T}D_{ii}R_{ij})_{pq}
=
\sum_{t=1}^{m}
(R_{ir})_{tp}(R_{ij})_{tq}d_{tt}
$$

Умножение на \(d_{tt}=\pm1\) не считаем

Для каждого \(t\) требуется одно умножение

$$
(R_{ir})_{tp}(R_{ij})_{tq}
$$

и одно вычитание полученного значения из элемента \(A_{rj}\)

Следовательно, для одного элемента требуется

$$
2m
$$

операций

В блоке \(m^2\) элементов, поэтому одно обновление блока требует

$$
2m\cdot m^2=2m^3
$$

операций

Теперь посчитаем число таких обновлений

Для каждого блока \(A_{ij}\), где

$$
2\le i\le j\le k,
$$

необходимо вычесть вклады всех предыдущих блочных строк:

$$
t=1,\ldots,i-1
$$

Таким образом, количество обновлений равно

$$
\sum_{i=2}^{k}
\sum_{j=i}^{k}
\sum_{t=1}^{i-1}1
$$

То есть

$$
\sum_{i=2}^{k}(k-i+1)(i-1)
$$

Сделаем замену \(s=i-1\):

$$
\sum_{s=1}^{k-1}s(k-s)
=
k\sum_{s=1}^{k-1}s
-
\sum_{s=1}^{k-1}s^2
$$

Отсюда

$$
\sum_{s=1}^{k-1}s(k-s)
=
k\frac{k(k-1)}2
-
\frac{k(k-1)(2k-1)}6
=
\frac{k(k-1)(k+1)}6
$$

Следовательно,

$$
C_3
=
2m^3
\frac{k(k-1)(k+1)}6
=
\frac{m^3k(k-1)(k+1)}3
$$

\subsubsection{Общая сложность построения \(R\) и \(D\)}

Складываем три полученных части:

$$
C_{\text{fact}}
=
C_1+C_2+C_3
$$

Получаем

$$
C_{\text{fact}}
=
\frac{k\,m(m+1)(2m+1)}6
+
\frac{k(k-1)m^3}{2}
+
\frac{k(k-1)(k+1)m^3}{3}
$$

Так как

$$
n=km,
$$

можно выразить результат через \(n\):

$$
C_{\text{fact}}
=
\frac{n^3}{3}
+
\frac{n^2m}{2}
-
\frac{nm^2}{2}
+
\frac{nm}{2}
+
\frac{n}{6}
$$

Или компактнее:

$$
C_{\text{fact}}
=
\frac{n^3}{3}
+
\frac{nm(n-m+1)}{2}
+
\frac{n}{6}
$$

Главный член при фиксированном размере блока \(m\):

$$
C_{\text{fact}}
\sim
\frac{n^3}{3}
$$

\subsection{Решение системы}

После получения разложения

$$
A=R^TDR
$$

решаем последовательно

$$
R^Tz=b,
$$

$$
Dy=z,
$$

$$
Rx=y
$$

\subsubsection{Решение \(R^Tz=b\)}

Внутри одного диагонального блока имеем

$$
(z_i)_r
=
\frac{
  (b_i)_r
  -
  \sum_{t=1}^{r-1}
  (R_{ii}^T)_{rt}(z_i)_t
}{
  (R_{ii}^T)_{rr}
}
$$

Для одного элемента требуется

$$
2(r-1)+1=2r-1
$$

операций

Поэтому решение одной треугольной системы размера \(m\) требует

$$
\sum_{r=1}^{m}(2r-1)=m^2
$$

операций

Таких диагональных блоков \(k\), следовательно,

$$
C_{\text{diag\,solve}}=km^2
$$

Перед решением очередного блока необходимо вычесть уже известные произведения

$$
b_i
-
\sum_{t=1}^{i-1}R_{ti}^{T}z_t
$$

Одно произведение матрицы размера \(m\times m\) на вектор с одновременным вычитанием требует

$$
2m^2
$$

операций

Число таких произведений равно

$$
\sum_{i=2}^{k}(i-1)
=
\frac{k(k-1)}2
$$

Поэтому

$$
C_{\text{off\,solve}}
=
2m^2\frac{k(k-1)}2
=
m^2k(k-1)
$$

Суммарно для первого шага:

$$
C_{R^Tz=b}
=
km^2+m^2k(k-1)
=
k^2m^2
$$

Так как \(n=km\),

$$
C_{R^Tz=b}=n^2
$$

\subsubsection{Вычисление \(Dy=z\)}

Матрица \(D\) диагональная и состоит из элементов \(\pm1\)

По принятому правилу умножение на \(\pm1\) не считаем, поэтому

$$
C_{Dy=z}=0
$$

\subsubsection{Решение \(Rx=y\)}

Вычисления полностью аналогичны первому шагу, только обход выполняется снизу вверх

Поэтому

$$
C_{Rx=y}=n^2
$$

Таким образом, общая стоимость поиска решения после разложения составляет

$$
C_{\text{solve}}
=
2n^2
$$

\subsection{Итоговая сложность}

Складываем стоимость разложения и решения:

$$
C
=
C_{\text{fact}}
+
C_{\text{solve}}
$$

Получаем

$$
C
=
\frac{n^3}{3}
+
\frac{n^2m}{2}
-
\frac{nm^2}{2}
+
\frac{nm}{2}
+
\frac{n}{6}
+
2n^2
$$

То есть

$$
\boxed{
  C
  =
  \frac{n^3}{3}
  +
  \left(\frac{m}{2}+2\right)n^2
  -
  \frac{m^2-m}{2}n
  +
  \frac{n}{6}
}
$$

При фиксированном размере блока \(m\):

$$
\boxed{
  C
  =
  \frac{n^3}{3}
  +
  O(n^2)
}
$$

и, следовательно,

$$
\boxed{
  C=O(n^3)
}
$$

\end{document}
